\section{Introduction}

Single-cell RNA sequencing (scRNA-seq) resolves gene expression at the level of
individual cells, enabling the study of cellular heterogeneity and the identification
of cell types~\citep{kiselev2019}. Unsupervised clustering is the primary route to cell
type identification, and it serves two distinct goals: recovering abundant cell types
and discovering rare populations whose low abundance makes them difficult to detect.
Both goals matter biologically, but they are not measured equally by the metrics in
common use.

Clustering methods for scRNA-seq fall into three families. The first reduces the
expression matrix and applies a conventional algorithm, for example principal component
analysis followed by $k$-means, or a shared-nearest-neighbor graph followed by Louvain
or Leiden community detection~\citep{leiden2019}. The second couples representation
learning with the clustering objective; deep embedded clustering (DEC)~\citep{dec2016}
and its single-cell descendants scDeepCluster~\citep{scdeepcluster2019} and
scDCC~\citep{scdcc2021} train an autoencoder while pulling the latent codes toward
trainable cluster centers. The third uses self-supervised objectives; scMAE~\citep{scmae2024}
perturbs gene expression, predicts which entries were corrupted, and reconstructs the
original values, learning representations that capture gene--gene dependencies.

A recent comprehensive benchmark reports a consistent and, at first sight, paradoxical
result: coupling the clustering objective into representation learning does not help.
Across a large and diverse collection of datasets, DEC ranks near the bottom of all
methods, while a self-supervised masked autoencoder without any clustering-specific loss
is among the most robust~\citep{scclubench2025}. The same benchmark attributes the weak
performance of most methods to a decoupling of embedding learning from clustering
optimization, which yields embedding spaces that are not conducive to clustering. This
raises a specific question: DEC does not decouple these objectives---it optimizes them
jointly---so why does it rank last?

We answer this question by opening the DEC latent space and measuring its per-dimension
variance. The confidence-gated Kullback--Leibler sharpening that defines DEC pulls cells
toward a small set of centers and, as a side effect, collapses the variance of most
latent coordinates: the effective dimensionality, measured as the participation ratio
of the per-dimension variance spectrum, falls from $114$ to $60$ on Macosko~\citep{macosko2015}.
Because $k$-means read-out operates on this collapsed space, its ability to separate fine
subpopulations degrades. The correction follows directly from the diagnosis. A
per-dimension variance floor, identical to the variance term of
VICReg~\citep{vicreg2022}, penalizes any coordinate whose batch standard deviation drops
below one and prevents the collapse. This choice is supported by two independent findings
from a self-supervised benchmark on single-cell data: VICReg-style variance regularization
is among the strongest generic mechanisms, and moderate to large embedding dimensionalities
consistently improve results~\citep{scsslbench2025}.

Two further design choices are grounded in the same benchmark and in our own ablations.
First, we retain zero-masking as the corruption operator rather than neighbor-based
augmentation, because random masking is reported to be the most effective augmentation
for single-cell data~\citep{scsslbench2025}, and because in our experiments neighbor
mixing has a net effect on ARI of $-0.001$, indistinguishable from zero. Second, we apply
the variance floor to the full clustering latent rather than to a reserved subspace,
because splitting the latent removes clustering capacity and collapses performance in
direct tests.

The contributions of this work are the following.
\begin{enumerate}
\item We diagnose per-dimension variance collapse as the failure mode of DEC on scRNA-seq
data, quantifying the drop in effective dimensionality from $114$ to $60$ and the
associated $2.5\times$ increase in seed-to-seed variance.
\item We correct it with a single VICReg variance-floor term inside an scMAE+DEC backbone,
improving ARI over DEC by $0.272$ on average (17 of 18 datasets) and restoring the
effective dimensionality to $114$.
\item We re-evaluate clustering under macro-F1 and rare-cluster recall, and show that ARI
systematically over-credits sharpening methods---VarFloor-scMAE moves from third to sixth
place when the ranking metric changes---and that rare-cell recall establishes no clear
winner among all methods examined.
\item We include a plain PCA+$k$-means baseline with a known number of clusters and full
dataset coverage, which proves unexpectedly strong and constrains the claims that can be
made for any deep method.
\end{enumerate}

We do not claim to surpass the scMAE backbone in accuracy. Our contribution is diagnostic
and corrective with respect to DEC, and evaluative with respect to rare-cell discovery.
